\chapter{Introduction and Architecture of Information Retrieval}

\section{Definition and Motivations}

\textbf{Information Retrieval} (IR) is the scientific and engineering discipline concerned with representing, storing, organizing, and accessing large collections of predominantly unstructured or semi-structured data to satisfy an \textbf{information need} expressed by a user~\cite{manning2008,buttcher2010}.

\begin{definition}[Information Retrieval (Manning et al.)]
Information Retrieval is finding material of an unstructured nature (usually text) that satisfies an information need from within large collections stored on computers.
\end{definition}

In the modern digital landscape, Information Retrieval extends far beyond general-purpose Web search engines (such as Google or Bing) to serve as the foundational infrastructure for:
\begin{itemize}
    \item \textbf{E-commerce and marketplaces:} product retrieval driven by textual descriptions, structured specifications, and customer reviews.
    \item \textbf{Enterprise and Desktop Search:} indexing and semantic searching across enterprise repositories (documents, emails, source code bases).
    \item \textbf{RAG (\textit{Retrieval-Augmented Generation}) Systems:} grounding Large Language Models (LLMs) with external, authoritative context documents to mitigate hallucinations and ensure up-to-date knowledge.
    \item \textbf{Autonomous agents and recommender engines:} content-based and collaborative filtering for personalized information discovery.
\end{itemize}

\subsection{Information Retrieval vs. Data Retrieval}
A fundamental conceptual distinction lies in contrasting traditional relational database querying (\textit{Data Retrieval}) with Information Retrieval:

\begin{table}[htbp]
\centering
\small
\begin{tabularx}{\textwidth}{p{3.2cm} X X}
\toprule
\textbf{Dimension} & \textbf{Data Retrieval (SQL Databases)} & \textbf{Information Retrieval} \\
\midrule
\textbf{Data Structure} & Highly structured data with strict relational schemas (tables, types, constraints) & Unstructured or semi-structured data (free text, HTML, PDF, metadata) \\
\addlinespace
\textbf{Query Language} & Formal deterministic queries (e.g., SQL) with exact relational algebra semantics & Natural language keywords or phrases (often ambiguous, short, and incomplete) \\
\addlinespace
\textbf{Matching Criterion} & Exact Boolean matching: a record either strictly satisfies or violates the predicate & Partial and probabilistic matching: documents exhibit varying degrees of relevance \\
\addlinespace
\textbf{Result Ordering} & No intrinsic ordering (unless explicitly specified via \texttt{ORDER BY}) & Central ranking mechanism driven by relevance estimation (\textit{Relevance Ranking}) \\
\addlinespace
\textbf{Error Tolerance} & Zero tolerance: syntax errors or unexpected values produce errors or empty sets & High tolerance: resilience to typos, synonyms, polysemy, and morphological variants \\
\bottomrule
\end{tabularx}
\caption{Analytical comparison between traditional Data Retrieval and Information Retrieval.}
\label{tab:ir-vs-dr-en}
\end{table}

\begin{alertblock}{Relevance: The Core of IR}
While database systems evaluate binary truth against relational predicates, Information Retrieval focuses on estimating the degree of \textbf{relevance} of each document with respect to the user's underlying cognitive information need. Relevance is inherently subjective, context-dependent, and time-varying.
\end{alertblock}

\section{The Two Fundamental Goals: Effectiveness and Efficiency}

A large-scale information retrieval system must simultaneously optimize two competing objectives:
\begin{enumerate}
    \item \textbf{Effectiveness:} Quantifies the intrinsic quality of the returned results, assessing how well the ranked list satisfies the user's search intent (thoroughly addressed in Chapter~\ref{chap:evaluation}).
    \item \textbf{Efficiency:} Quantifies computational cost, throughput (queries processed per second), and latency (demanded within $10$--$50$ milliseconds per query across billions of documents).
\end{enumerate}

\subsection{Efficiency and Underlying Hardware Architecture}
The efficiency of a modern search engine depends not merely on algorithmic complexity, but on hardware execution dynamics:
\begin{itemize}
    \item \textbf{Pipelining and Superscalar Processors:} Modern CPUs segment instruction execution into sequential stages (Instruction Fetch, Decode, Execute, Memory Access, Write Back), issuing multiple instructions per clock cycle.
    \item \textbf{Pipeline Hazards:}
    \begin{itemize}
        \item \textit{Structural Hazards:} Resource contention when multiple instructions compete for the same execution unit or memory bus.
        \item \textit{Data Hazards:} Logical data dependencies between consecutive instructions (RAW: Read-After-Write), forcing pipeline stalls (\textit{bubbles}) when values are unavailable.
        \item \textit{Control Hazards:} Branch mispredictions. During posting list intersection loops, an incorrect branch prediction flushes the pipeline, incurring a penalty of $15$--$20$ clock cycles per misprediction.
    \end{itemize}
    \item \textbf{Memory Hierarchy:} CPU registers operate at sub-nanosecond speeds, L1/L2/L3 caches respond in a few clock cycles, whereas main memory (RAM) requires $50$--$100$ ns and secondary storage (NVMe/SSD) demands tens of microseconds. Posting list layouts prioritize sequential, contiguous memory layouts to maximize cache-line reuse and trigger hardware prefetching.
    \item \textbf{SIMD (\textit{Single Instruction Multiple Data}) Vectorization:} Utilizing 128-bit, 256-bit, or 512-bit vector registers (AVX-2, AVX-512, ARM NEON) enables simultaneous comparisons of 8 or 16 document identifiers in a single machine cycle, dramatically accelerating posting intersections.
    \item \textbf{Parallel Computing and Distributed Sharding:} Indices exceeding the capacity of a single machine are partitioned (\textit{sharded}), typically through document-level sharding, allowing compute clusters to process queries concurrently with centralized result aggregation.
\end{itemize}

\section{Information Retrieval System Architecture}

As depicted in Figure~\ref{fig:ir-architecture-detailed-en}, an IR system enforces a strict decoupling between the \textbf{Offline Pipeline} (indexing and model training) and the \textbf{Online Pipeline} (low-latency query serving).

\begin{figure}[htbp]
\centering
\begin{tikzpicture}[
    scale=0.85,
    box/.style={rectangle, rounded corners=3pt, draw=blue!75!black, fill=blue!8, thick, minimum width=2.5cm, minimum height=0.75cm, text width=2.4cm, align=center, font=\scriptsize},
    qbox/.style={rectangle, rounded corners=3pt, draw=blue!75!black, fill=blue!15, thick, minimum width=2.1cm, minimum height=0.65cm, text width=2.0cm, align=center, font=\scriptsize},
    db/.style={cylinder, shape border rotate=90, aspect=0.22, draw=cyan!75!black, fill=cyan!10, thick, minimum width=2.3cm, minimum height=0.9cm, text width=2.2cm, inner sep=2pt, align=center, font=\scriptsize\bfseries},
    serpbox/.style={rectangle, rounded corners=3pt, draw=teal!80!black, fill=teal!12, thick, minimum width=2.1cm, minimum height=0.75cm, text width=2.0cm, align=center, font=\scriptsize},
    arr/.style={-Stealth, thick, blue!80!black}
]
    % Dotted Boundary Line (Online / Offline)
    \draw[blue!40!black, dashed, thick] (-3.6, 2.8) -- (11.6, 2.8);
    \node[draw=blue!50!black, fill=blue!5, rounded corners=2pt, font=\bfseries\scriptsize, text=blue!80!black, inner sep=3pt, anchor=east] at (-2.0, 3.4) {ONLINE};
    \node[draw=blue!50!black, fill=blue!5, rounded corners=2pt, font=\bfseries\scriptsize, text=blue!80!black, inner sep=3pt, anchor=east] at (-2.0, 2.2) {OFFLINE};

    % ONLINE LAYER
    % Column 1: Query Pipeline
    \node[qbox] (query) at (0, 7.6) {Query};
    \node[qbox] (expquery) at (0, 6.3) {Expanded Query};
    \node[box] (qproc) at (0, 5.0) {Query Processing};

    % Column 2: Feature Lookup
    \node[box] (featlook) at (3.6, 5.0) {Feature Lookup\\and Computation};

    % Column 3: Learned Ranking
    \node[box] (rankfunc) at (7.2, 5.0) {Learned Ranking\\Function};
    \node[serpbox] (serp) at (10.4, 5.0) {\textbf{Ranked SERP}\\(Results)};

    % SHARED / BOUNDARY STORAGE (Cylinders)
    \node[db] (invidx) at (0, 2.8) {Inverted\\Index};
    \node[db] (docfeat) at (3.6, 2.8) {Document Features\\Repository};
    \node[db] (ltrmodel) at (7.2, 2.8) {Learning-to-Rank\\Model};

    % OFFLINE LAYER
    % Middle tier: Processors
    \node[box] (indexing) at (0, 0.8) {Indexing};
    \node[box] (featproc) at (3.6, 0.8) {Feature Processor};
    \node[box] (training) at (7.2, 0.8) {Training};

    % Bottom tier: Sources (Cylinders)
    \node[db] (doccoll) at (0, -1.2) {Document\\Collection};
    \node[db] (traindata) at (7.2, -1.2) {Training\\Data};

    % ARROWS - ONLINE
    \draw[arr] (query) -- (expquery);
    \draw[arr] (expquery) -- (qproc);
    \draw[arr] (invidx) -- (qproc);
    \draw[arr] (qproc) -- (featlook);
    \draw[arr] (docfeat) -- (featlook);
    \draw[arr] (featlook) -- (rankfunc);
    \draw[arr] (ltrmodel) -- (rankfunc);
    \draw[arr] (rankfunc) -- (serp);

    % ARROWS - OFFLINE
    \draw[arr] (doccoll) -- (indexing);
    \draw[arr] (indexing) -- (invidx);
    \draw[arr] (doccoll.north east) -- (featproc.south west);
    \draw[arr] (featproc) -- (docfeat);
    \draw[arr] (traindata) -- (training);
    \draw[arr] (training) -- (ltrmodel);
\end{tikzpicture}
\caption{Grid architecture of a modern search engine (Online vs. Offline pipelines).}
\label{fig:ir-architecture-detailed-en}
\end{figure}

\subsection{Text Processing Pipeline}
Raw textual content undergoes four sequential preparation stages:
\begin{enumerate}
    \item \textbf{Tokenization:} Decomposing character streams into lexical tokens, handling punctuation, special characters, and hyphenated terms.
    \item \textbf{Normalization and Case Folding:} Converting tokens to lowercase and stripping accents/diacritics to ensure uniformity across queries and documents.
    \item \textbf{Stopword Filtering:} Removing high-frequency functional words with minimal discriminative value (e.g., \texttt{"the"}, \texttt{"is"}, \texttt{"at"}, \texttt{"which"}).
    \item \textbf{Stemming and Lemmatization:}
    \begin{itemize}
        \item \textit{Stemming:} Algorithmic stripping of suffixes using heuristic rules (such as Porter's algorithm~\cite{manning2008}) to extract the root form (\textit{stem}). Example: \texttt{"retrieving"}, \texttt{"retrieval"} $\to$ \texttt{"retriev"}.
        \item \textit{Lemmatization:} Full grammatical analysis utilizing vocabulary dictionaries to map inflected forms to their formal base lemma. Example: \texttt{"am"}, \texttt{"were"} $\to$ \texttt{"be"}.
    \end{itemize}
\end{enumerate}

\section{Language Properties and Retrieval Models}

Natural language word distributions adhere to statistical empirical laws that shape index designs and ranking algorithms:
\begin{itemize}
    \item \textbf{Zipf's Law:} The frequency $f(r)$ of a term is inversely proportional to its frequency rank $r$:
    \begin{equation}
    f(r) \propto \frac{1}{r}
    \end{equation}
    A handful of words dominate overall occurrences, while the vast majority of words appear extremely rarely (\textit{long tail}).
    \item \textbf{Heaps' Law:} The vocabulary size $|V|$ grows as a power law with respect to collection size in terms of total word tokens $N$:
    \begin{equation}
    |V| = k \cdot N^\beta \quad (\text{with } 0.4 \le \beta \le 0.6)
    \end{equation}
    The vocabulary never fully saturates as more documents are gathered.
\end{itemize}

Core retrieval models include:
\begin{enumerate}
    \item \textbf{Boolean Model:} Deterministic matching against logical operators ($\land, \lor, \neg$), yielding unranked result sets.
    \item \textbf{Vector Space Model (TF-IDF):} Documents and queries are modeled as vectors across vocabulary space $|V|$, weighted by term frequency ($TF$) and inverse document frequency ($IDF$), scored via cosine similarity.
    \item \textbf{Probabilistic Model (BM25):} The industry standard for lexical retrieval based on the 2-Poisson indexing model, incorporating non-linear term frequency saturation and document length normalization~\cite{buttcher2010}.
\end{enumerate}

\section{The Inverted Index}

The \textbf{Inverted Index} reverses document storage directionality, mapping terms to posting lists $t \mapsto \mathcal{P}_t$~\cite{manning2008}.

\begin{definition}[Inverted Index]
An Inverted Index consists of two complementary structures:
\begin{enumerate}
    \item \textbf{Dictionary (Vocabulary):} Stores all unique terms appearing across the collection, maintained via fast lookup structures (B-Trees or Hash Tables), along with document frequency $df_t$ and a pointer to the start of its posting list.
    \item \textbf{Posting List:} For each term $t$, an ordered array of document identifiers sorted strictly in ascending order ($\text{DocID}_1 < \text{DocID}_2 < \dots$). Each posting may also record local term frequencies ($tf$) and intra-document term position offsets.
\end{enumerate}
\end{definition}

\begin{figure}[htbp]
\centering
\begin{tikzpicture}[
    scale=0.95,
    term/.style={rectangle, draw=blue!80!black, fill=blue!10, thick, minimum width=2.4cm, minimum height=0.68cm, font=\small\ttfamily, align=left},
    df/.style={rectangle, draw=blue!80!black, fill=blue!5, thick, minimum width=0.9cm, minimum height=0.68cm, font=\small, align=center},
    post/.style={rectangle, draw=purple!80!black, fill=purple!8, thick, minimum width=0.85cm, minimum height=0.68cm, font=\small\bfseries},
    ptr/.style={-Stealth, thick, blue!80!black},
    link/.style={-Stealth, thick, purple!80!black},
    header/.style={font=\bfseries\small, align=center}
]
    % Main Headers
    \node[header, blue!80!black] at (0.825, 4.0) {DICTIONARY};
    \node[header, purple!80!black] at (5.3, 4.0) {POSTING LISTS};

    % Sub-headers
    \node[font=\footnotesize\itshape, blue!60!black] at (0, 3.4) {Term};
    \node[font=\footnotesize\itshape, blue!60!black] at (1.65, 3.4) {$df_t$};
    \node[font=\footnotesize\itshape, purple!70!black] at (5.3, 3.4) {$\mathrm{DocID}_1 < \mathrm{DocID}_2 < \dots$};

    % Dictionary column (Terms + Doc Frequency df_t)
    \node[term] (t1) at (0, 2.7) {an};
    \node[df]   (df1) at (1.65, 2.7) {2};

    \node[term] (t2) at (0, 1.8) {inverted};
    \node[df]   (df2) at (1.65, 1.8) {3};

    \node[term] (t3) at (0, 0.9) {index};
    \node[df]   (df3) at (1.65, 0.9) {3};

    \node[term] (t4) at (0, 0.0) {query};
    \node[df]   (df4) at (1.65, 0.0) {1};

    % Postings
    \node[post] (p11) at (4.2, 2.7) {1};
    \node[post] (p12) at (5.3, 2.7) {2};

    % Row 2
    \node[post] (p21) at (4.2, 1.8) {1};
    \node[post] (p22) at (5.3, 1.8) {2};
    \node[post] (p23) at (6.4, 1.8) {3};

    % Row 3
    \node[post] (p31) at (4.2, 0.9) {1};
    \node[post] (p32) at (5.3, 0.9) {2};
    \node[post] (p33) at (6.4, 0.9) {3};

    % Row 4
    \node[post] (p41) at (4.2, 0.0) {3};

    % Arrows from Dictionary to Posting Lists
    \draw[ptr] (df1.east) -- (p11.west);
    \draw[ptr] (df2.east) -- (p21.west);
    \draw[ptr] (df3.east) -- (p31.west);
    \draw[ptr] (df4.east) -- (p41.west);

    % Links between postings
    \draw[link] (p11) -- (p12);
    \draw[link] (p21) -- (p22);
    \draw[link] (p22) -- (p23);
    \draw[link] (p31) -- (p32);
    \draw[link] (p32) -- (p33);
\end{tikzpicture}
\caption{Inverted Index structure: Term dictionary with document frequency $df_t$ and linked posting lists ordered by ascending DocID.}
\label{fig:inverted-index-structure-en}
\end{figure}

\subsection{Query Processing and Merge Intersection}
Maintaining strictly ascending DocID order allows conjunctive (\textbf{AND}) queries to be resolved via a synchronized two-pointer merge algorithm (Algorithm~\ref{alg:merge-intersection-en}).

\begin{algorithm}[H]
\caption{Merge Intersection for Conjunctive Queries ($P_1 \land P_2$)}
\label{alg:merge-intersection-en}
\KwInput{Two sorted posting lists $P_1$ and $P_2$}
\KwOutput{Sorted list $R$ containing DocIDs present in both lists}
$R \leftarrow \emptyset$\;
$i \leftarrow 1$; $j \leftarrow 1$\;
\While{$i \le |P_1|$ \textbf{and} $j \le |P_2|$}{
    \If{$P_1[i] == P_2[j]$}{
        $R.\text{append}(P_1[i])$\;
        $i \leftarrow i + 1$\;
        $j \leftarrow j + 1$\;
    }
    \ElseIf{$P_1[i] < P_2[j]$}{
        $i \leftarrow i + 1$\;
    }
    \Else{
        $j \leftarrow j + 1$\;
    }
}
\Return{$R$}\;
\end{algorithm}

The worst-case runtime complexity is $\BigO{|P_1| + |P_2|}$. Disjunctive queries (\textbf{OR}) compute the union of the posting lists, while Phrase Search checks positional difference constraints across word occurrences.

\section{Index Compression Techniques}

Compressing inverted indices serves two goals: minimizing storage footprint and converting slow I/O bus transfer bottlenecks into fast CPU decompression operations.

\subsection{Binary Interpolative Coding}
\textbf{Binary Interpolative Coding} (introduced by Moffat and Stuiver~\cite{moffat2000}) is a specialized compression algorithm tailored for strictly monotonically increasing integer sequences.

\begin{definition}[Interpolative Coding Principle]
Given a sorted sub-sequence of $n$ integers $L = [x_1, x_2, \dots, x_n]$ bounded in the interval $[l, h]$ (such that $l \le x_1 < x_2 < \dots < x_n \le h$), the algorithm selects the median element $m = \lfloor (n+1)/2 \rfloor$.
Because there must be at least $m - 1$ distinct integers preceding $x_m$ and at least $n - m$ distinct integers following it, $x_m$ is restricted to the interval:
\begin{equation}
l + m - 1 \le x_m \le h - n + m
\end{equation}
The number of valid candidate values for $x_m$ is:
\begin{equation}
R = (h - n + m) - (l + m - 1) + 1 = h - l - n + 1
\end{equation}
The shifted value $(x_m - (l + m - 1))$ is encoded using exactly $\lceil \log_2 R \rceil$ bits. The procedure then recurses on the left sub-array over $[l, x_m - 1]$ and on the right sub-array over $[x_m + 1, h]$.
\end{definition}

\begin{example}[Recursive Decomposition Example from Slides]
Consider the sequence of $n = 12$ identifiers within $[0, 62]$:
\[
L = [3, 4, 7, 13, 14, 15, 21, 25, 36, 38, 54, 62] \quad (n=12, m=6, l=0, h=62)
\]
The recursive tree partitions the sequence as follows:
\begin{itemize}
    \item The median splits the array into left sublist $[3, 4, 7, 13, 14]$ with $(n=5, m=3, l=0, h=14)$ and right sublist $[21, 25, 36, 38, 54]$ with $(n=5, m=3, l=16, h=62)$.
    \item At the next depth, $[3, 4, 7, 13, 14]$ splits into $[3, 4]$ with $(n=2, m=1, l=0, h=6)$ and $[13, 14]$ with $(n=2, m=1, l=8, h=14)$.
    \item The leaves encode individual elements: for instance, $[14]$ with $(n=1, m=1, l=14, h=14)$ yields $R = 14 - 14 - 1 + 1 = 1$, requiring $\lceil \log_2 1 \rceil = 0$ bits!
\end{itemize}
When postings are densely clustered, $R$ collapses to $1$, yielding zero-bit representations for clustered entries.
\end{example}

\subsection{Dictionary-Based Compression: LZ77}
For general text or token streams, the \textbf{Lempel-Ziv '77} (LZ77) algorithm maintains a sliding window comprising a Search Buffer and a Lookahead Buffer, emitting triples:
\begin{equation}
(\text{offset}, \text{length}, \text{next\_char})
\end{equation}
For example, given the input sequence:
\[
T = \text{\texttt{a b a b a a a a b a b a a a a b}}
\]
LZ77 generates the pointer sequence:
\[
(0, \text{\texttt{'a'}}), \; (0, \text{\texttt{'b'}}), \; (2, 3), \; (1, 3), \; (7, 8)
\]
replacing repeated substrings with compact backward references.

\section{Evolution Toward Neural Information Retrieval}

Beginning in 2018, Transformer architectures and pretrained models like BERT revolutionized Information Retrieval~\cite{khattab2020,formal2021}. Neural architectures split into two primary paradigms:

\begin{figure}[htbp]
\centering
\begin{tikzpicture}[
    box/.style={rectangle, rounded corners=4pt, draw=blue!80!black, fill=blue!8, thick, minimum width=2.4cm, minimum height=0.8cm, inner sep=4pt, align=center, font=\small},
    proc/.style={rectangle, rounded corners=4pt, draw=teal!80!black, fill=teal!8, thick, minimum width=2.5cm, minimum height=1.15cm, inner sep=6pt, align=center, font=\small},
    score/.style={rectangle, rounded corners=4pt, draw=purple!80!black, fill=purple!10, thick, minimum width=2.8cm, minimum height=0.85cm, inner sep=5pt, align=center, font=\small\bfseries},
    arr/.style={-Stealth, thick, blue!80!black},
    title/.style={font=\bfseries\normalsize, align=center}
]
    % Divider
    \draw[gray!35, dashed, thick] (6.6, 5.0) -- (6.6, -0.5);

    % --- INTERACTION-BASED (CROSS-ENCODER) ---
    \node[title, blue!80!black] at (3.0, 4.4) {Cross-Encoder\\[2pt]\footnotesize (Interaction-based)};

    \node[box] (q1) at (1.6, 3.2) {Query $q$};
    \node[box] (d1) at (4.4, 3.2) {Document $d$};
    \node[proc, minimum width=5.2cm] (cross) at (3.0, 1.7) {Cross-Encoder (BERT)\\[3pt]\footnotesize Full Cross-Attention over all tokens};
    \node[score] (s1) at (3.0, 0.1) {Score $s(q, d)$};

    \draw[arr] (q1.south) -- (cross.140);
    \draw[arr] (d1.south) -- (cross.40);
    \draw[arr] (cross.south) -- (s1.north);

    % --- REPRESENTATION-BASED (BI-ENCODER) ---
    \node[title, teal!80!black] at (10.2, 4.4) {Bi-Encoder\\[2pt]\footnotesize (Representation-based)};

    \node[box] (q2) at (8.8, 3.2) {Query $q$};
    \node[box] (d2) at (11.6, 3.2) {Document $d$};
    \node[proc, minimum width=2.5cm] (encq) at (8.8, 1.7) {Query Encoder\\[3pt]\footnotesize $\vec{e}_q = f_\theta(q)$};
    \node[proc, minimum width=2.5cm] (encd) at (11.6, 1.7) {Doc Encoder\\[3pt]\footnotesize $\vec{e}_d = g_\theta(d)$};
    \node[score] (s2) at (10.2, 0.1) {Similarity $\vec{e}_q \cdot \vec{e}_d$};

    \draw[arr] (q2.south) -- (encq.north);
    \draw[arr] (d2.south) -- (encd.north);
    \draw[arr] (encq.south) -- (s2.140);
    \draw[arr] (encd.south) -- (s2.40);
\end{tikzpicture}
\caption{Structural comparison between Cross-Encoder (Interaction-based) and Bi-Encoder (Representation-based).}
\label{fig:cross-vs-bi-encoder-en}
\end{figure}

\begin{enumerate}
    \item \textbf{Interaction-Based Models (Cross-Encoders):} The query and document are concatenated into a joint token stream:
    \[
    \text{Input} = \texttt{[CLS]} \circ q \circ \texttt{[SEP]} \circ d
    \]
    and fed together through all self-attention layers. This captures fine-grained cross-token interactions, delivering top-tier ranking quality. However, its high inference cost ($\BigO{|\mathcal{D}| \times \text{Transformer}}$) limits its application to second-stage re-ranking over candidate pools (e.g., top-100 results).
    \item \textbf{Representation-Based Models (Bi-Encoders):} Query and document are encoded independently into vectors. Document embeddings can be indexed offline, reducing online matching to inner products or vector search.
\end{enumerate}

\subsection{Three Embedding Families}
Representation models fall into three broad categories:
\begin{enumerate}
    \item \textbf{Dense Single-Vector (e.g., DPR, Contriever):} Each document and query is mapped to a single dense vector in $\mathbb{R}^d$ ($d \approx 768$), queried via graph-based vector indices (like HNSW).
    \item \textbf{Dense Multi-Vector (e.g., ColBERT~\cite{khattab2020}):} Generates token-level embeddings for query and document tokens, evaluated via \textbf{Late Interaction} using the \textit{MaxSim} operator:
    \begin{equation}
    S(q, d) = \sum_{t_q \in q} \max_{t_d \in d} \left( \vec{E}_{t_q} \cdot \vec{E}_{t_d}^T \right)
    \end{equation}
    Preserves token identity and achieves superior generalization, at the cost of larger index sizes.
    \item \textbf{Sparse Single-Vector (e.g., SPLADE~\cite{formal2021}):} Encodes queries and documents into high-dimensional sparse vectors over the vocabulary $|V|$ ($|V| \approx 30\,000$). Non-zero values represent learned term weights including automated term expansions.
    \begin{itemize}
        \item \textbf{Effectiveness:} High performance across both in-domain and zero-shot out-of-domain evaluation.
        \item \textbf{Efficiency:} Highly compact, directly indexable inside standard Inverted Index structures.
        \item \textbf{Interpretability:} Weights align with human-readable vocabulary words, enabling transparent debugging.
    \end{itemize}
\end{enumerate}

\section{\texorpdfstring{$k$-Approximate Nearest Neighbor Search ($k$-ANN)}{k-Approximate Nearest Neighbor Search (k-ANN)}}

In high-dimensional embedding spaces, retrieval reduces to finding the top-$k$ documents maximizing the inner product with the query vector (\textit{Maximum Inner Product Search}, MIPS):
\begin{equation}
\operatorname{argmax}_{v \in \mathcal{D}}^{(k)} \; \vec{q}^{\,T} \cdot \vec{v}
\end{equation}
On standard Web-scale collections:
\begin{itemize}
    \item The document volume exceeds $8.8$ million items (e.g., MS MARCO benchmark).
    \item The vocabulary vector space comprises roughly $30\,000$ dimensions.
    \item Each document activates only approximately $150$ non-zero components.
\end{itemize}
High-throughput query serving avoids brute-force linear scans through hybrid forward-inverted sparse vector index architectures, exemplified by modern libraries such as \textbf{TusKANNy}, engineered in Rust at the University of Pisa.
